\section{Noncompact integration of the radial derivative terms}
\label{rf-section}

The remaining radial integration takes place on the whole complex
plane.  Both the point $t=0$ and infinity must therefore be treated.
We prove the Fourier regularization and the actual elementary
differential identities used in this passage
\cite{alpo-uniform-source}.  Two simplifications are useful.
Only the Fourier transform of the weighted density is needed.
Furthermore, for an admissible word the differential expansion
can be written explicitly as a sum over subsets.
The estimates for the actual Jacquet integrals needed to apply
these facts are identified separately at the end.

\subsection{Fourier normalization and extension across zero}

We identify $\mathbb C$ with $\mathbb R^2$ and use the measure
$dt=dx\,dy$ for $t=x+iy$.  The Fourier transform in this section is
\begin{equation}\label{rf-transform}
 \widehat E(s)=\int_{\mathbb C}
        \exp(st-\overline s\,\overline t)\,E(t)\,dt.
\end{equation}
For $s=a+ib$, the exponential in \eqref{rf-transform} is
$\exp(2i(bx+ay))$.  In particular its absolute value is one.
All derivatives below are real derivatives.  The norm
$\|D^jF(t)\|$ is the operator norm of the $j$-linear derivative
on the Euclidean space $\mathbb R^2$; for $j=0$ it is $|F(t)|$.

\begin{lemma}\label{rf-inversion}
Let $E\in L^1(\mathbb C)$ and suppose
$\widehat E\in L^1(\mathbb C)$.  Then
\begin{equation}\label{rf-inverse}
 E_0(t)=\frac1{\pi^2}\int_{\mathbb C}
       \exp(-st+\overline s\,\overline t)\widehat E(s)\,ds
\end{equation}
is continuous and equals $E$ almost everywhere.
It equals $E$ at every point at which $E$ is continuous.
If $|s|^j\widehat E(s)$ is integrable for $0\leq j\leq m$,
then $E_0$ is $C^m$, and differentiation under the integral is valid.
\end{lemma}

\begin{proof}
The absolute convergence and continuity of \eqref{rf-inverse}
follow from dominated convergence.  To fix the normalization
and prove the assertion about $E$, put
\[
 \eta_a(t)=\frac1{\pi a}e^{-|t|^2/a}\qquad(a>0).
\]
The one-dimensional Gaussian integral, applied separately to
the two real coordinates of $s$, gives
\[
 \widehat{\eta_a}(s)=e^{-a|s|^2},\qquad
 \frac1{\pi^2}\int_{\mathbb C}e^{-a|s|^2}
        e^{st-\overline s\,\overline t}\,ds=\eta_a(t).
\]
Fubini's theorem is applicable because $E$ is integrable and
the Gaussian in $s$ is integrable.  It therefore gives
\[
 (E*\eta_a)(t)=\frac1{\pi^2}\int_{\mathbb C}
      e^{-st+\overline s\,\overline t}
      e^{-a|s|^2}\widehat E(s)\,ds.
\]
The right side tends to $E_0(t)$ uniformly in $t$ as $a\to0$.

On the other hand, $E*\eta_a\to E$ in $L^1$.
Indeed,
\[
 \|E*\eta_a-E\|_1
 \leq\int_{\mathbb C}\eta_a(v)\,
                    \|E(\,\cdot-v)-E\|_1\,dv.
\]
Translations are continuous in $L^1$: this follows first for
continuous functions of compact support by uniform continuity,
and then for an arbitrary integrable function by approximation
in $L^1$.  Split the last integral into a fixed small disk
and its complement, use this translation continuity on the
disk and the bound $2\|E\|_1$ on the complement.
The Gaussian mass of the complement tends to zero.
An $L^1$-convergent subsequence converges almost everywhere,
so the uniform limit $E_0$ equals $E$ almost everywhere.
At a point of continuity of $E$, their difference is continuous.
If its value were nonzero, its absolute value would be bounded
away from zero throughout a small disk, contradicting equality
almost everywhere.

For a unit real direction $v$, differentiating the exponential
in \eqref{rf-inverse} gives the factor
$-sv+\overline s\,\overline v$, whose absolute value is at most
$2|s|$.  The stated moment assumptions permit repeated
differentiation by dominated convergence and prove the last
assertion.
\end{proof}

\begin{theorem}\label{rf-weighted-extension}
Let $m\geq0$ be an integer and let $D:\mathbb C\to\mathbb C$
be measurable and continuous on $\mathbb C^\times$.
Put
\[
                   E(t)=(1+|t|^2)^2D(t).
\]
Suppose that $E\in L^1(\mathbb C)$ and, for some $C\geq0$,
\begin{equation}\label{rf-weighted-decay}
             |\widehat E(s)|\leq C(1+|s|)^{-m-3}
                      \qquad(s\in\mathbb C).
\end{equation}
Then there is a unique $C^m$ function $D_0$ on $\mathbb C$
equal to $D$ on $\mathbb C^\times$.  It satisfies
\begin{align}
 \int_{\mathbb C}D_0(t)\,dt
      &=\int_{\mathbb C}D(t)\,dt, \label{rf-same-integral}\\
 \|D^jD_0(t)\|&\leq A_m C(1+|t|)^{-4}
       \qquad(0\leq j\leq m), \label{rf-all-derivatives}\\
 \int_{\mathbb C}D^jD_0(t)\,dt&=0
       \qquad(1\leq j\leq m). \label{rf-zero-derivatives}
\end{align}
In the last line the integral is the componentwise integral
of the multilinear derivative.  The constant $A_m$ depends
only on $m$, and the following completely explicit choice is valid:
\begin{align*}
 B_0&=1,\qquad B_j=\prod_{\nu=0}^{j-1}(5\nu+4),\\
 M_{m,j}&=\frac{2^{j+1}}{\pi}
                   \frac{(j+1)!(m-j)!}{(m+2)!}
                       \qquad(0\leq j\leq m),\\
 A_m&=\max_{0\leq j\leq m}
       \sum_{\nu=0}^j\binom j\nu
                    16\,2^{\nu/2}B_\nu M_{m,j-\nu}.
\end{align*}
No separate estimate for $\widehat D$ is required.
\end{theorem}

\begin{proof}
The hypothesis $E\in L^1$ implies $D\in L^1$ because
$(1+|t|^2)^{-2}\leq1$.  For $0\leq j\leq m$, polar coordinates
and the substitution $r=u/(1-u)$ give
\begin{align*}
 \int_{\mathbb C}|s|^j(1+|s|)^{-m-3}\,ds
 &=2\pi\int_0^\infty r^{j+1}(1+r)^{-m-3}\,dr\\
 &=2\pi\frac{(j+1)!(m-j)!}{(m+2)!}.
\end{align*}
The last integral can also be evaluated by repeated integration
by parts; its endpoint terms vanish since $j\leq m$.
Lemma~\ref{rf-inversion} supplies a $C^m$ extension $E_0$ and gives
\begin{equation}\label{rf-E-derivative}
                       \|D^jE_0(t)\|\leq M_{m,j}C.
\end{equation}
Since $E$ is continuous away from zero, $E_0=E$ there.
Define $D_0(t)=(1+|t|^2)^{-2}E_0(t)$.  This proves the asserted
extension and uniqueness, because $\mathbb C^\times$ is dense.

For clarity, a uniform derivative bound for the weight can
be obtained without an unspecified differentiation constant.
For every coordinate multi-index $\alpha$ with $|\alpha|=j$,
\[
 \partial^\alpha(1+x^2+y^2)^{-2}
     =\frac{P_\alpha(x,y)}{(1+x^2+y^2)^{2+j}},
 \qquad \deg P_\alpha\leq j.
\]
The sum of the absolute values of the coefficients of
$P_\alpha$ is at most $B_j$.  To see this, differentiating once
replaces $P$ by
\[
  (1+x^2+y^2)\partial_iP-(4+2j)x_iP.
\]
The coefficient sum increases by a factor at most
$3j+4+2j=5j+4$.  This proves the bound by induction.
On $|t|\leq1$ and $|t|\geq1$, respectively, the displayed
rational expression then gives
\[
 |\partial^\alpha(1+|t|^2)^{-2}|
                      \leq16B_j(1+|t|)^{-4}.
\]
Passing from coordinate derivatives to the operator norm
costs at most $2^{j/2}$, since
$|v_x|+|v_y|\leq\sqrt2\,|v|$.
The multilinear product rule, together with
\eqref{rf-E-derivative}, proves \eqref{rf-all-derivatives}
with the stated $A_m$.

These bounds imply integrability of $D_0$ and all of its
derivatives through order $m$.  Equality almost everywhere
with $D$ proves \eqref{rf-same-integral}.
To prove \eqref{rf-zero-derivatives}, it is enough to consider
a coordinate derivative $\partial_i H$, where
$H=\partial^\alpha D_0$ has order at most $m-1$.
Choose a smooth compactly supported function $\chi$ equal to
one on the unit disk and put $\chi_R(t)=\chi(t/R)$.
Integration by parts on its compact support gives
\[
 \int_{\mathbb C}\chi_R\,\partial_i H
        =-\int_{\mathbb C}H\,\partial_i\chi_R.
\]
The right side has absolute value at most
$R^{-1}\|\partial_i\chi\|_\infty\|H\|_1$, hence tends to zero.
Dominated convergence on the left gives
$\int\partial_iH=0$.  The argument applies to every component
of every positive-order derivative.  It treats zero and
infinity simultaneously, without discarding an annular
boundary term.
\end{proof}

\begin{corollary}\label{rf-finite-combination}
Let $D_1,\ldots,D_N$ satisfy
Theorem~\ref{rf-weighted-extension} with a common order $m$,
and constants $C_i$ that may depend on $i$.  Let $0\leq j_i\leq m$,
$a_i\in\mathbb C$ and fixed real directions
$v_{i,1},\ldots,v_{i,j_i}\in\mathbb C$.
If, for $t\ne0$,
\[
 f(t)=\sum_{i=1}^N a_i
      D^{j_i}D_i(t)[v_{i,1},\ldots,v_{i,j_i}],
\]
then $f$ extends continuously to $\mathbb C$, its extension
is $O((1+|t|)^{-4})$, and
\[
        \int_{\mathbb C}f(t)\,dt
           =\sum_{\{i:j_i=0\}}a_i\int_{\mathbb C}D_i(t)\,dt.
\]
The implicit constant can be taken to be
$A_m\sum_i |a_i|C_i\prod_{\nu=1}^{j_i}|v_{i,\nu}|$.
The assertion includes $N=0$.
\end{corollary}

\begin{proof}
Replace $D_i$ by its extension from the theorem.
They agree on an open neighborhood of every nonzero point,
so their derivatives agree there.  The finite sum of the
extended derivatives is continuous and has the stated
bound.  Integrate term by term and use
\eqref{rf-same-integral}--\eqref{rf-zero-derivatives}.
\end{proof}

\begin{example}\label{rf-critical-power}
The extra integrable Fourier moment cannot simply be omitted.
For $t\ne0$ set
\[
 E(t)=\int_0^1\frac1{\pi a}e^{-|t|^2/a}\,da,\qquad
 D(t)=\frac{E(t)}{(1+|t|^2)^2},
\]
and give both functions any finite values at zero.
Tonelli's theorem gives $\|E\|_1=1$, and differentiation
under the integral shows that $E$ is smooth away from zero.
Fubini's theorem and the Gaussian calculation above give
\[
 \widehat E(s)=
 \begin{cases}
 (1-e^{-|s|^2})/|s|^2,&s\ne0,\\
 1,&s=0.
 \end{cases}
\]
Thus $|\widehat E(s)|\leq4(1+|s|)^{-2}$.
For $0<|t|<1$, restricting the integral to
$|t|^2\leq a\leq1$ gives
$E(t)\geq(\pi e)^{-1}\log(|t|^{-2})$.
Consequently $D$ has no continuous extension at zero.
This is a counterexample to replacing the power $-3$
by $-2$ in the case $m=0$; the source uses the sufficient
power $-m-3$.
\end{example}

\subsection{The actual two-component differential identity}

Let $\mathcal H$ be a complex Hilbert space, with inner product
linear in the first argument.  Write
\[
 \partial_t=\tfrac12(\partial_x-i\partial_y),\qquad
 \partial_{\overline t}=\tfrac12(\partial_x+i\partial_y),
 \qquad p=t\partial_t,\quad
 \overline p=\overline t\,\partial_{\overline t}.
\]
Here $\overline p$ denotes this displayed differential
operator, not an instruction to conjugate its argument.
For $u,v:\mathbb C^\times\to\mathcal H^2$ put
\[
 D(u,v)(t)=|t|^{-2}
              \bigl(\langle u_0(t),v_0(t)\rangle+
                    \langle u_1(t),v_1(t)\rangle\bigr).
\]
For a real number $Y$, define
\begin{align}
 V_+(Y)v&=(Yv_0+pv_0-tv_1,\ t^{-1}v_0),\nonumber\\
 I_+(Y)u&=(-\overline t^{-1}u_1,\
             \overline t\,u_0+(Y+1)u_1-\overline p\,u_1),
                    \label{rf-plus-operators}\\
 V_-(Y)v&=(Yv_0-\overline p\,v_0-\overline t\,v_1,\
                                      \overline t^{-1}v_0),\nonumber\\
 I_-(Y)u&=(-t^{-1}u_1,\
                    t u_0+(Y-1)u_1+p u_1).
                    \label{rf-minus-operators}
\end{align}
These are the source's actual one-coordinate operators, including
their shifts by $1$.  Let $Ev=(0,v_0)$.

\begin{lemma}\label{rf-one-letter}
For $C^1$ Hilbert-space-valued functions $u,v$ on
$\mathbb C^\times$ and real $Y$,
\begin{align}
 D(I_+(Y)u,V_+(Y)v)
      &=D(u,v)-\partial_{\overline t}D(u,Ev),\label{rf-plus-density}\\
 D(I_-(Y)u,V_-(Y)v)
      &=D(u,v)+\partial_tD(u,Ev).\label{rf-minus-density}
\end{align}
\end{lemma}

\begin{proof}
Put $A(t)=\langle u_1(t),v_0(t)\rangle$.
Direct expansion of \eqref{rf-plus-operators}, using
linearity in the first argument and $Y=\overline Y$, gives
\[
 \langle I_+u,V_+v\rangle
 =\langle u,v\rangle+\overline t^{-1}A
       -\langle\overline t^{-1}\overline p\,u_1,v_0\rangle
       -\langle u_1,t^{-1}p v_0\rangle.
\]
The last two terms sum to $-\partial_{\overline t}A$:
\[
 \partial_{\overline t}A
    =\langle\partial_{\overline t}u_1,v_0\rangle
       +\langle u_1,\partial_t v_0\rangle.
\]
Since
$\partial_{\overline t}|t|^{-2}
       =-\overline t^{-1}|t|^{-2}$,
multiplication by $|t|^{-2}$ proves \eqref{rf-plus-density}.

The analogous expansion of \eqref{rf-minus-operators} is
\[
 \langle I_-u,V_-v\rangle
 =\langle u,v\rangle-t^{-1}A
       +\langle t^{-1}p u_1,v_0\rangle
       +\langle u_1,\overline t^{-1}\overline p\,v_0\rangle.
\]
The last two terms sum to $\partial_tA$.
Using $\partial_t|t|^{-2}=-t^{-1}|t|^{-2}$ proves
\eqref{rf-minus-density}.  All calculations are valid in a
Hilbert space by the differentiability of its inner product.
\end{proof}

\subsection{An explicit expansion for an admissible word}

Let $S$ be a finite set and let
$\mathcal B=\{0,1\}^S$.  Replace $\mathcal H^2$ by
$\mathcal H^{\mathcal B}$ and sum its inner product over
$\mathcal B$.  For $a\in S$, each of the operators
\eqref{rf-plus-operators}--\eqref{rf-minus-operators} acts
on the $a$-th two-component factor and as the identity
on the other factors.  Denote by $E_a$ the corresponding
elementary matrix $E_{10}$:
\[
 (E_a v)_b=
 \begin{cases}
 v_{b[a\leftarrow0]},& b(a)=1,\\
 0,&b(a)=0.
 \end{cases}
\]

A word of length $h$ is a sequence
$\ell_j=(a_j,\sigma_j,Y_j)$, where
$a_j\in S$, $\sigma_j\in\{+,-\}$ and $Y_j\in\mathbb R$.
Assume that the $a_j$ are pairwise distinct.
This is the part of the source's admissibility condition
used in the present identity.  Define the ordered products
\[
 I_w=I_{\sigma_1,a_1}(Y_1)\cdots I_{\sigma_h,a_h}(Y_h),
 \qquad
 V_w=V_{\sigma_1,a_1}(Y_1)\cdots V_{\sigma_h,a_h}(Y_h).
\]
The rightmost factor acts first.  Put
\[
 (\epsilon_j,\delta_j)=
 \begin{cases}
 (-1,\partial_{\overline t}),&\sigma_j=+,\\
 (1,\partial_t),&\sigma_j=-.
 \end{cases}
\]
For $J\subseteq\{1,\ldots,h\}$, let
$E_J=\prod_{j\in J}E_{a_j}$,
$\epsilon_J=\prod_{j\in J}\epsilon_j$, and
$\delta_J=\prod_{j\in J}\delta_j$.
The factors in these three products commute.
For the empty set they are the identity, $1$, and the
identity differential operator, respectively.

\begin{proposition}\label{rf-subset-expansion}
For smooth $u,v:\mathbb C^\times\to\mathcal H^{\mathcal B}$,
\begin{equation}\label{rf-explicit-expansion}
             D(I_wu,V_wv)
       =\sum_{J\subseteq\{1,\ldots,h\}}
                         \epsilon_J\delta_JD(u,E_Jv).
\end{equation}
There are at most $2^h$ terms, their coefficients have
absolute value one, and the derivative order of the
$J$-term is $|J|$.
\end{proposition}

\begin{proof}
The case $h=0$ is immediate.  Write $w=\ell_1w'$.
Sum Lemma~\ref{rf-one-letter} over the other component
indices to obtain
\[
 D(I_wu,V_wv)=D(I_{w'}u,V_{w'}v)
     +\epsilon_1\delta_1
          D(I_{w'}u,E_{a_1}V_{w'}v).
\]
Since $a_1$ does not occur in $w'$, $E_{a_1}$ commutes
with every constant component matrix in $V_{w'}$.
It also commutes with multiplication by $t,\overline t$
and with their derivatives.  Hence
$E_{a_1}V_{w'}=V_{w'}E_{a_1}$.
Apply the induction hypothesis first to $(u,v)$ and
then to $(u,E_{a_1}v)$.  The resulting terms correspond,
respectively, to subsets not containing $1$ and to
subsets containing $1$.  Smoothness permits the
indicated differentiations of these identities.
\end{proof}

\begin{theorem}\label{rf-noncompact-step}
In Proposition~\ref{rf-subset-expansion}, suppose for every
$J\subseteq\{1,\ldots,h\}$ that
\[
 E^{(J)}(t)=(1+|t|^2)^2D(u,E_Jv)(t)
\]
is integrable and
\[
             |\widehat {E^{(J)}}(s)|
                    \leq C_J(1+|s|)^{-h-3}
                    \qquad(s\in\mathbb C)
\]
for some $C_J\geq0$.
Then $D(I_wu,V_wv)$ has a continuous integrable extension
to the whole plane, bounded in absolute value by
$A_h(\sum_J C_J)(1+|t|)^{-4}$, and
\begin{equation}\label{rf-full-plane-identity}
       \int_{\mathbb C}D(I_wu,V_wv)(t)\,dt
                       =\int_{\mathbb C}D(u,v)(t)\,dt.
\end{equation}
No support condition in $t$ is imposed.
\end{theorem}

\begin{proof}
Apply Theorem~\ref{rf-weighted-extension} to the finitely
many functions $D(u,E_Jv)$, with $m=h$.
An iterated Wirtinger derivative of order $j$ is a sum
of $2^j$ coordinate derivatives, each with coefficient
of absolute value $2^{-j}$.  Thus \eqref{rf-all-derivatives}
bounds it by the same constant, and its integral is
zero for $j\geq1$ by \eqref{rf-zero-derivatives}.
Insert these extended functions in
\eqref{rf-explicit-expansion}.  The empty subset contributes
$D(u,v)$; every other subset integrates to zero.
\end{proof}

\begin{example}\label{rf-repeated-index}
The distinct-index assumption in
Proposition~\ref{rf-subset-expansion} is essential.
Take $\mathcal H=\mathbb C$, one two-component factor,
two identical plus letters with $Y=0$, and
$u=(0,1)$, $v=(1,0)$.
Then $I_+^2u=(-\overline t^{-1},0)$ and
$V_+^2v=(-1,0)$, so
\[
                D(I_+^2u,V_+^2v)=\frac1{t\overline t^2}.
\]
The proposed subset expression would instead be
$-2\partial_{\overline t}|t|^{-2}
       =2/(t\overline t^2)$, since $E^2=0$.
The source's admissible words have distinct indices,
so this does not contradict its statement.
\end{example}

\subsection{Relation to the radial form and remaining analysis}

The measure $|t|^{-2}dt$ explains the shifts in the
operators.  On compactly supported smooth functions
away from zero, integration by parts gives
\[
       p^*=-\overline p,\qquad
       \overline p^{\,*}=-p,\qquad t^*=\overline t.
\]
These formulas define the formal adjoint of every
matrix differential operator considered here.

\begin{lemma}\label{rf-inverse-adjoint}
As differential operators on $\mathbb C^\times$,
\[
             I_\sigma(Y)V_\sigma(Y)^*
             =V_\sigma(Y)^*I_\sigma(Y)=1
                      \qquad(Y\in\mathbb R).
\]
Consequently $I_wV_w^*=1$ for an ordered word.
If matrix differential operators $R,R'$ satisfy
$R=V_w^*R'V_w$, then
\begin{equation}\label{rf-mirror-identity}
                         R'V_w=I_wR.
\end{equation}
\end{lemma}

\begin{proof}
In the plus case,
\[
 V_+(Y)^*=
  \begin{pmatrix}Y-\overline p&\overline t^{-1}\\
                          -\overline t&0\end{pmatrix},
 \qquad
 I_+(Y)=
  \begin{pmatrix}0&-\overline t^{-1}\\
                  \overline t&Y+1-\overline p\end{pmatrix}.
\]
Their products are the identity because
$[\overline p,\overline t]=\overline t$ and
$[\overline p,\overline t^{-1}]=-\overline t^{-1}$.
In the minus case the corresponding matrices are
\[
 V_-(Y)^*=
  \begin{pmatrix}Y+p&t^{-1}\\-t&0\end{pmatrix},
 \qquad
 I_-(Y)=
  \begin{pmatrix}0&-t^{-1}\\t&Y-1+p\end{pmatrix},
\]
and the calculation uses $[p,t]=t$ and
$[p,t^{-1}]=-t^{-1}$.
In the word product $I_wV_w^*$ these identities cancel
adjacent factors, starting with the last letter.
Multiplying $R=V_w^*R'V_w$ on the left by $I_w$
proves \eqref{rf-mirror-identity}.
These are algebraic identities; they assert no
unproved adjoint-domain statement on a noncompact space.
\end{proof}

\begin{lemma}\label{rf-L2-smoothness}
Let $(Z,\nu)$ be a $\sigma$-finite measure space and let
$F:\mathbb C^\times\times Z\to\mathbb C^d$ be measurable.
Suppose that, outside one null subset of $Z$, it is smooth
in $t$ and all its coordinate derivatives are jointly
measurable.  Suppose also that, for every compact
$K\subset\mathbb C^\times$ and every multi-index $\alpha$,
\begin{equation}\label{rf-local-L2-bounds}
       \sup_{t\in K}
          \|\partial^\alpha_tF(t,\cdot)\|_{L^2(Z,\nu)}
                                                     <\infty.
\end{equation}
Then $t\mapsto F(t,\cdot)$ is a smooth map into
$L^2(Z,\nu;\mathbb C^d)$, with the expected derivatives.
A pointwise dominating $L^2$ function is not required.
\end{lemma}

\begin{proof}
Fix a closed disk $K$ inside $\mathbb C^\times$ and a
smaller concentric disk.  Every sufficiently short
segment from the smaller disk lies in $K$.
The fundamental theorem of calculus, applied outside
the common null set, gives
\[
 F(t+h,z)-F(t,z)
       =\int_0^1\bigl(h_x\partial_xF+h_y\partial_yF\bigr)
                                                    (t+uh,z)\,du.
\]
Cauchy--Schwarz in $u$, followed by Tonelli's theorem,
or equivalently the integral triangle inequality in
$L^2$, bounds the $L^2$ norm of the difference by
$\sqrt2|h|$ times a common bound for the two first
derivatives on $K$.  This proves norm continuity.
Applying the same argument to every
$\partial^\alpha F$ proves their norm continuity.

The second-order Taylor formula along that segment
gives
\[
 F(t+h,\cdot)-F(t,\cdot)
       -h_x\partial_xF(t,\cdot)-h_y\partial_yF(t,\cdot)
 =\int_0^1(1-u)
       \sum_{i,j\in\{x,y\}}h_i h_j\partial_{ij}F(t+uh,\cdot)\,du.
\]
If $B$ bounds the $L^2$ norms of the second derivatives
on $K$, the norm of the right side is at most
$B(|h_x|+|h_y|)^2/2\leq B|h|^2$.
Thus the derivative exists in $L^2$ and has the asserted
value.  Repeating this proof for each
$\partial^\alpha F$ proves smoothness of all orders.
The same estimates justify every integral displayed
above; no assumption of finite measure on $Z$ is used.
\end{proof}

For the source radial form, $\mathcal H$ is the $L^2$ space
of the remaining radial coordinates and the terminal
unitary group, with its specified positive measure
and Hilbert--Schmidt coefficient space.  Its measure
is independent of the coordinate $t$ being integrated.
When the radial functions and their derivatives are
smooth as maps into this Hilbert space, all the preceding
identities apply directly.  Lemma~\ref{rf-L2-smoothness}
shows that it suffices to bound their $L^2$ derivative
norms locally uniformly.  The source's coordinate
operators include multiplication by $t^{\pm1}$ and
$\overline t^{\pm1}$ together with $p,\overline p$.
On a compact subset of $\mathbb C^\times$, every
coordinate derivative of order $j$ is a finite sum
of words in these operators, of length at most $2j$.
For example $\partial_t=t^{-1}p$ and
$\partial_{\overline t}=\overline t^{-1}\overline p$,
and one repeats these identities by composition.
Consequently locally uniform $L^2$ bounds for all
coordinate-operator words imply
\eqref{rf-local-L2-bounds}.  This also justifies
differentiating all the integrated pairings in
Lemma~\ref{rf-one-letter}, without constructing a
separate pointwise dominating function.

With $u=Rx$ and $v=y$, identity
\eqref{rf-mirror-identity} and Theorem~\ref{rf-noncompact-step}
give
\[
 \int_{\mathbb C}D(R'V_wx,V_wy)(t)\,dt
                  =\int_{\mathbb C}D(Rx,y)(t)\,dt,
\]
provided the displayed weighted Fourier bounds hold for
the densities $D(Rx,E_Jy)$.
This is precisely the one-coordinate integration step
after the source's algebraic relation
$R=V_w^*R'V_w$ has been established.
The argument uses the actual shifts and the full
noncompact coordinate domain.

The present section proves the elementary differential
expansion, the inverse formal-adjoint identity, and the
analytic removal of every derivative term.  It does
not yet prove that all actual Jacquet-generated functions
satisfy the required $L^2$ differentiability and weighted
Fourier estimates, nor does it identify the preceding
algebraic relation with all of the source's recursively
defined frame operators.  The source obtains the
Fourier estimates from four specified differential
operators preserving the class of coefficients and
from an identity that transfers multiplication by the
Fourier parameter to those operators.  Their construction
and the precise Jacquet pairing and measure identities
remain substantive dependencies.  In particular,
pointwise terminal positivity alone is not being used
as an integrated noncompact theorem.


